% Part: first-order-logic
% Chapter: axiomatic-deduction
% Section: proof-theoretic-notions

\documentclass[../../../include/open-logic-section]{subfiles}

\begin{document}

\iftag{FOL}
      {\olfileid[hi]{fol}{axd}{ptn}}
      {\olfileid[hi]{pl}{axd}{ptn}}

\olsection{प्रमाण-सैद्धांतिक धारणाएँ}

\begin{explain}
जिस प्रकार हमने कई महत्त्वपूर्ण अर्थवैज्ञानिक धारणाएँ
(\iftag{FOL}{वैधता}{सर्वसत्यता}, तार्किक निष्पत्ति, संतुष्टनीयता) परिभाषित
की हैं, उसी प्रकार अब उनकी संगत \emph{प्रमाण-सैद्धांतिक धारणाएँ} परिभाषित
करते हैं। इनकी परिभाषा में !!{sentence}s की !!{structure}s में संतुष्टि का
सहारा नहीं लिया जाता, बल्कि कुछ सूत्रों की !!{derivability} या
!!{nonderivability} का सहारा लिया जाता है। इन धारणाओं का एक-दूसरे से मेल
खाना एक महत्त्वपूर्ण खोज थी। यह मेल \emph{सुदृढ़ता} और \emph{पूर्णता
प्रमेयों} का विषय है।
\end{explain}

\begin{defn}[!!^{derivability}]
!!^a{formula}~$!A$, $\Gamma$ से \emph{!!{derivable}} है, जिसे
$\Gamma \Proves !A$ लिखते हैं, यदि~$\Gamma$ से ऐसी !!a{derivation} हो जो
~$!A$ पर समाप्त होती है।
\end{defn}

\begin{defn}[प्रमेय]
!!^a{formula}~$!A$ \emph{प्रमेय} है यदि रिक्त समुच्चय से $!A$ की कोई
!!a{derivation} हो। हम $\Proves !A$ लिखते हैं यदि $!A$ प्रमेय है, और यदि
नहीं है तो
$\Proves/ !A$ लिखते हैं।
\end{defn}

\begin{defn}[संगतता]
!!{formula}s का समुच्चय $\Gamma$ \emph{संगत} तभी और केवल तभी है जब
$\Gamma\Proves/ \lfalse$; अन्यथा वह \emph{असंगत} है।
\end{defn}

\begin{prop}[स्वतुल्यता]
\ollabel{prop:reflexivity}
यदि $!A \in \Gamma$, तो $\Gamma \Proves !A$।
\end{prop}

\begin{proof}
  अकेला !!{formula}~$!A$ ही~$!A$ की,~$\Gamma$ से, !!a{derivation} है।
\end{proof}

\begin{prop}[एकदिष्टता]
\ollabel{prop:monotonicity}
यदि $\Gamma \subseteq \Delta$ और $\Gamma \Proves !A$, तो $\Delta
\Proves !A$।
\end{prop}

\begin{proof}
$!A$ की, $\Gamma$ से, प्रत्येक !!{derivation}, $!A$ की,~$\Delta$ से,
!!a{derivation} भी है।
\end{proof}

\begin{prop}[संक्रामकता]
\ollabel{prop:transitivity}
यदि $\Gamma \Proves !A$ और $\{!A\} \cup \Delta \Proves
!B$, तो $\Gamma \cup \Delta \Proves !B$।
\end{prop}

\begin{proof}
  मान लें $\{!A\} \cup \Delta \Proves !B$। तब कोई !!a{derivation}
  $!B_1$, \dots, $!B_l = !B$,~$\{!A\} \cup
  \Delta$ से, है। उस
  !!{derivation} के कुछ चरण सही होंगे क्योंकि कोई नियम पहले की
  पंक्ति~$!B_i = !A$ का निर्देश करता है। परिकल्पना से~$!A$ की,~$\Gamma$ से,
  कोई !!a{derivation} है, अर्थात् !!a{derivation}~$!A_1$, \dots,
  $!A_k = !A$, जिसमें प्रत्येक $!A_i$ कोई अभिगृहीत,~$\Gamma$ का
  !!a{element}, अथवा
  किसी अनुमान-नियम से सही चरण है। अब इस अनुक्रम पर विचार करें:
  \[
  !A_1, \dots, !A_k = !A, !B_1, \dots, !B_l = !B.
  \]
  यह~$!B$ की, $\Gamma \cup \Delta$ से, सही !!{derivation} है, क्योंकि अब हर
  $B_i = !A$ को वही नियम उचित ठहराता है जो~$!A_k = !A$ को उचित ठहराता है।
\end{proof}

ध्यान दें कि विशेषतः इसका अर्थ है: यदि $\Gamma \Proves !A$ और $!A
\Proves !B$, तो $\Gamma \Proves !B$। इससे यह भी मिलता है कि यदि $!A_1,
\dots, !A_n \Proves !B$ और $\Gamma \Proves !A_i$ प्रत्येक~$i$ के लिए हो,
तो $\Gamma \Proves !B$।

\begin{prop}
\ollabel{prop:incons}
$\Gamma$ असंगत है तभी और केवल तभी जब $\Gamma \Proves !A$, प्रत्येक~$!A$
के लिए।
\end{prop}

\begin{proof}
अभ्यास।
\end{proof}

\tagprob{FOL}
\begin{prob}
\olref[fol][axd][ptn]{prop:incons} सिद्ध कीजिए।
\end{prob}
\tagendprob

\tagprob{notFOL}
\begin{prob}
\olref[pl][axd][ptn]{prop:incons} सिद्ध कीजिए।
\end{prob}
\tagendprob

\begin{prop}[संहतता]
\ollabel{prop:proves-compact}
  \begin{enumerate}
  \item यदि $\Gamma \Proves !A$, तो कोई परिमित उपसमुच्चय $\Gamma_0
    \subseteq \Gamma$ ऐसा है कि $\Gamma_0 \Proves !A$।
  \item यदि~$\Gamma$ का हर परिमित उपसमुच्चय संगत है, तो $\Gamma$ संगत है।
  \end{enumerate}
\end{prop}

\begin{proof}
  \begin{enumerate}
    \item यदि $\Gamma \Proves !A$, तो !!{formula}s का कोई परिमित अनुक्रम
      $!A_1$, \dots,~$!A_n$ ऐसा है कि $!A \ident !A_n$, और प्रत्येक
      $!A_i$ या तो कोई तार्किक अभिगृहीत,~$\Gamma$ का !!a{element}, अथवा
      पूर्वपक्ष-स्थापन द्वारा पिछले !!{formula}s से प्राप्त है। $\Gamma_0$
      को उन $!A_i$ का समुच्चय मानें जो~$\Gamma$ में हैं। तब वही
      !!{derivation},~$\Gamma_0$ से भी !!a{derivation} है; अतः
      $\Gamma_0 \Proves !A$।
    \item यह विशेष प्रसंग $!A
      \ident \lfalse$ के लिए~(1) का प्रतिधन कथन
      है।
\end{enumerate}
\end{proof}

\end{document}
